% ============================================================================
%  videotest: учебное пособие по проекту
%  Компиляция: xelatex tutorial.tex   (рекомендуется)
%              pdflatex tutorial.tex  (работает, если есть babel-russian)
%  Нужны пакеты: babel, fontspec+polyglossia (xelatex) или babel-russian
%  (pdflatex), amsmath, amssymb, mathtools, booktabs, longtable, listings,
%  hyperref, xcolor, geometry, titlesec, enumitem, microtype.
%  Шрифты CMU (для xelatex): texlive-fonts-extra или пакет cm-unicode.
% ============================================================================
\documentclass[11pt,a4paper]{article}

% --- кодировка и язык -------------------------------------------------------
\usepackage{iftex}
\ifPDFTeX
  \usepackage[T2A]{fontenc}
  \usepackage[utf8]{inputenc}
  \usepackage[russian,english]{babel}
\else
  \usepackage{fontspec}
  \usepackage{polyglossia}
  \setdefaultlanguage{russian}
  \setotherlanguage{english}
  \setmainfont{CMU Serif}
  \setsansfont{CMU Sans Serif}
  \setmonofont{CMU Typewriter Text}
\fi

% --- математика -------------------------------------------------------------
\usepackage{amsmath,amssymb,amsthm,mathtools}
\usepackage{bm}

% --- вёрстка ----------------------------------------------------------------
\usepackage[a4paper,margin=2.4cm]{geometry}
\usepackage{booktabs}
\usepackage{array}
\usepackage{longtable}
\usepackage{enumitem}
\usepackage{xcolor}
\usepackage{graphicx}
\usepackage{caption}
\usepackage{titlesec}
\usepackage{microtype}

% --- код --------------------------------------------------------------------
\usepackage{listings}
\definecolor{codebg}{gray}{0.96}
\definecolor{codeframe}{gray}{0.75}
\definecolor{codekw}{RGB}{0,70,140}
\definecolor{codecomment}{RGB}{0,110,60}
\definecolor{codestring}{RGB}{150,50,20}

\lstdefinestyle{cpp}{
  language=C++,
  basicstyle=\ttfamily\footnotesize,
  keywordstyle=\color{codekw}\bfseries,
  commentstyle=\color{codecomment}\itshape,
  stringstyle=\color{codestring},
  backgroundcolor=\color{codebg},
  frame=single,
  rulecolor=\color{codeframe},
  breaklines=true,
  breakatwhitespace=false,
  postbreak=\mbox{\textcolor{red}{$\hookrightarrow$}\space},
  showstringspaces=false,
  tabsize=4,
  columns=fullflexible,
  keepspaces=true,
  extendedchars=true,
  inputencoding=utf8,
  literate={--}{{-{}-}}1 {<<}{{<}{<}}1 {>>}{{>}{>}}1,
}
\lstdefinestyle{shell}{
  language=bash,
  basicstyle=\ttfamily\footnotesize,
  backgroundcolor=\color{codebg},
  frame=single,
  rulecolor=\color{codeframe},
  breaklines=true,
  showstringspaces=false,
  columns=fullflexible,
}
\lstdefinestyle{cmake}{
  basicstyle=\ttfamily\footnotesize,
  backgroundcolor=\color{codebg},
  frame=single,
  rulecolor=\color{codeframe},
  breaklines=true,
  showstringspaces=false,
  columns=fullflexible,
}
\lstset{style=cpp}

% --- гиперссылки ------------------------------------------------------------
\usepackage[hidelinks,bookmarksnumbered]{hyperref}

% --- теоремы и примечания ---------------------------------------------------
\theoremstyle{definition}
\newtheorem{definition}{Определение}[section]
\newtheorem{example}[definition]{Пример}
\theoremstyle{remark}
\newtheorem*{remark}{Замечание}

% «Заметка» как среда: не требует tcolorbox и допускает абзацы внутри.
\newenvironment{note}
  {\par\medskip\noindent
   \begin{minipage}{\linewidth}
   \colorbox{codebg}{\begin{minipage}{\dimexpr\linewidth-2\fboxsep\relax}
     \small\textbf{Замечание.}\par\vspace{1pt}\noindent\ignorespaces}
  {\end{minipage}}\end{minipage}\par\medskip}

\titleformat{\section}{\Large\bfseries}{\thesection}{1em}{}
\titleformat{\subsection}{\large\bfseries}{\thesubsection}{1em}{}

\setlist[itemize]{topsep=2pt,itemsep=1pt,parsep=0pt}
\setlist[enumerate]{topsep=2pt,itemsep=1pt,parsep=0pt}

% ============================================================================
\title{\textbf{videotest}\\[2mm]
       \large Учебное пособие: архитектура C++-проекта, обработка видео\\
       и математика компьютерного зрения}
\author{Учебные материалы по проекту \texttt{videotest}}
\date{}

\begin{document}
\maketitle

\begin{abstract}
\noindent
Пособие разбирает учебный проект \texttt{videotest}: захват видео с камеры
через V4L2, инференс ONNX-модели instance-сегментации, сопровождение объектов
фильтром Калмана и одновременный показ исходного и размеченного потоков.
Документ рассчитан на читателя с сильной математической подготовкой и опытом
в нейросетях, но давно не писавшего на C++. Поэтому язык C++ объясняется
подробно: владение памятью, времена жизни объектов, многопоточность,
система сборки. Отдельные главы посвящены математике проекта: преобразованию
координат letterbox, декодированию выходов YOLO, подавлению пересечений (NMS),
построению масок из прототипов, фильтру Калмана и метрикам качества.

Всё изложенное соответствует коду в каталоге \texttt{src/}; ссылки на файлы
и функции даны по ходу текста.
\end{abstract}

\tableofcontents
\newpage

% ============================================================================
\section{Введение: что делает программа}

\subsection{Постановка задачи}

Программа получает поток кадров с веб-камеры (\texttt{/dev/videoX}), для
каждого кадра запускает нейросетевую модель instance-сегментации, сопоставляет
найденные объекты между кадрами (трекинг) и выводит на экран два изображения
рядом: слева --- исходный кадр, справа --- тот же кадр с наложенными масками,
рамками и идентификаторами объектов.

\begin{center}
\begin{tabular}{ll}
\toprule
Источник кадров & V4L2 (Video4Linux2), прямой доступ к драйверу \\
Модель & ONNX Runtime, YOLO-подобная сеть instance-сегментации \\
Трекинг & фильтр Калмана + сопоставление по IoU \\
Вывод & OpenCV HighGUI, два потока рядом \\
Потоки & 2 рабочих + главный (отображение) \\
\bottomrule
\end{tabular}
\end{center}

\subsection{Почему именно такая архитектура}

Три ограничения задают всё остальное.

\begin{enumerate}
  \item \textbf{Камера диктует темп.} При 30 кадрах в секунду новый кадр
        появляется каждые $33{,}3$~мс, и драйвер имеет конечное число
        внутренних буферов. Если их вовремя не возвращать, камера
        останавливается.
  \item \textbf{Инференс непредсказуем по времени.} Один кадр может
        обрабатываться 8~мс, а может 40~мс --- в зависимости от содержимого,
        загрузки CPU и модели.
  \item \textbf{Экран должен показывать свежее.} Если обработка отстаёт,
        разумно пропустить устаревшие кадры, а не выстраивать их в очередь.
\end{enumerate}

Отсюда всё устройство: камера и модель работают в разных потоках, между ними
находится очередь \emph{ограниченной ёмкости}, которая при переполнении
выбрасывает самый старый кадр. Задержка при этом остаётся постоянной, а не
растёт линейно со временем.

\subsection{Быстрый старт}

\begin{lstlisting}[style=shell,escapeinside={(*@}{@)}]
sudo apt install build-essential cmake libopencv-dev v4l-utils   (*@\#@*) OpenCV 4.5+
./scripts/fetch_onnxruntime.sh                                   (*@\#@*) SDK в third_party/
cmake -S . -B build -DCMAKE_BUILD_TYPE=Release
cmake --build build -j
ctest --test-dir build --output-on-failure
./build/videotest --model models/yolo11n-seg.onnx --classes 80
\end{lstlisting}

Без камеры и без весов модели:

\begin{lstlisting}[style=shell]
./scripts/smoke_test.sh
./build/videotest --synthetic --no-window --max-frames 60 \
                  --model build/stub-raw.onnx --classes 3 --snapshot out.png
\end{lstlisting}

\subsection{Структура каталогов}

\begin{lstlisting}[style=shell]
videotest/
  CMakeLists.txt          описание сборки
  README.md               краткая инструкция
  docs/                   документация
    architecture.md       механика проекта (кратко)
    tutorial.tex          это пособие
  scripts/
    fetch_onnxruntime.sh  скачивание SDK ONNX Runtime
    make_stub_model.py    генерация модели-заглушки для тестов
    smoke_test.sh         сквозная проверка без камеры и весов
  src/
    frame.hpp             структуры данных, ходящие между потоками
    bounded_queue.hpp     очередь с вытеснением
    logger.{hpp,cpp}      потокобезопасные логи
    stop_flag.{hpp,cpp}   кооперативная остановка по Ctrl+C
    main.cpp              разбор аргументов, цикл отображения
    streams/
      source.hpp          абстрактный интерфейс источника
      v4l2_source.*       камера через V4L2
      synthetic_source.*  синтетический источник (тесты)
    model/
      segmenter.*         ONNX: препроцессинг, декодирование, маски
      class_names.*       имена классов (COCO или файл)
    track/
      tracker.*           фильтр Калмана + IoU-сопоставление
    pipeline.*            два потока и очередь между ними
    render/
      renderer.*          отрисовка масок, рамок, HUD, склейка
  tests/
    queue_test.cpp
    tracker_test.cpp
\end{lstlisting}

Разделение на каталоги отражает \emph{направление зависимостей}: \texttt{streams}
и \texttt{model} не знают друг о друге, \texttt{pipeline} знает об обоих,
\texttt{main} собирает всё вместе. Такое разложение позволяет тестировать
трекер вообще без камеры и без модели.

% ============================================================================
\section{Как устроен C++-проект и как он собирается}

Этот раздел --- для читателя, который помнит язык, но забыл инструменты.

\subsection{Единицы трансляции, заголовки, линковка}

Программа на C++ не компилируется целиком. Каждый файл \texttt{.cpp}
(\emph{единица трансляции}) компилируется независимо в объектный файл
\texttt{.o}; затем \emph{компоновщик} (linker) соединяет объектные файлы и
библиотеки в исполняемый файл.

\begin{itemize}
  \item \textbf{Заголовочный файл} (\texttt{.hpp}) содержит \emph{объявления}:
        типы, сигнатуры функций, шаблоны, встраиваемые (\texttt{inline})
        функции. Он вставляется в \texttt{.cpp} директивой
        \texttt{\#include} на этапе препроцессора.
  \item \textbf{Файл реализации} (\texttt{.cpp}) содержит \emph{определения}:
        тела функций, глобальные объекты.
  \item Директива \texttt{\#pragma once} (или include guard) защищает от
        повторного включения одного заголовка в одной единице трансляции.
\end{itemize}

\begin{remark}
Шаблоны классов и функций определены \emph{полностью} в заголовках. Причина:
компилятор должен видеть тело шаблона, чтобы сгенерировать код для конкретного
типа аргумента. Именно поэтому \texttt{Queue<T>} в проекте --- один заголовок
\texttt{bounded\_queue.hpp} без \texttt{.cpp}.
\end{remark}

\subsection{Статическая и динамическая линковка}

\begin{center}
\begin{tabular}{p{0.28\linewidth}p{0.62\linewidth}}
\toprule
Вид & Что происходит \\
\midrule
Статическая (\texttt{.a}) & Код библиотеки \emph{копируется} в исполняемый файл
при компоновке. Файл больше, но зависимостей во время работы нет. \\
Динамическая (\texttt{.so}) & В исполняемый файл попадает только имя символа;
код подгружается при запуске. Обновление библиотеки не требует пересборки. \\
\bottomrule
\end{tabular}
\end{center}

В проекте используется оба вида: собственная библиотека
\texttt{videotest\_core} статическая (её код всё равно нужен целиком), а
OpenCV и ONNX Runtime --- динамические.

\begin{note}
Динамическая библиотека ищется во время выполнения, а не компиляции.
Если \texttt{libonnxruntime.so} лежит в нестандартном каталоге, программа
запустится с ошибкой <<cannot open shared object file>>. Решение ---
переменная \texttt{LD\_LIBRARY\_PATH} либо \emph{rpath}, вшитый в бинарник при
сборке (см.~\texttt{BUILD\_RPATH} в \texttt{CMakeLists.txt}).
\end{note}

\subsection{Система сборки CMake}

CMake не компилирует сам: он \emph{генерирует} файлы сборки (Makefile или
\texttt{build.ninja}) для выбранного генератора. Процесс двухфазный:

\begin{lstlisting}[style=shell,escapeinside={(*@}{@)}]
cmake -S . -B build -DCMAKE_BUILD_TYPE=Release   (*@\#@*) конфигурация
cmake --build build -j                            (*@\#@*) сборка
\end{lstlisting}

Ключевые понятия:

\begin{itemize}
  \item \textbf{Цель (target)} --- то, что собирается: исполняемый файл,
        библиотека или тест. Цель хранит свои свойства: исходники, флаги,
        зависимости.
  \item \textbf{Область видимости} \texttt{PRIVATE}/\texttt{PUBLIC}/\texttt{INTERFACE}.
        \texttt{PUBLIC} означает <<мне нужно и всем, кто меня линкует>>.
        В проекте \texttt{videotest\_core} объявляет зависимости
        (\texttt{OpenCV}, \texttt{ONNXRUNTIME}) как \texttt{PUBLIC} ---
        поэтому \texttt{videotest} и тесты получают их автоматически.
  \item \textbf{Тип сборки} (\texttt{CMAKE\_BUILD\_TYPE}) задаёт набор флагов:
        \texttt{Debug} ($-O0 -g$), \texttt{Release} ($-O3 -DNDEBUG$),
        \texttt{RelWithDebInfo} ($-O2 -g$).
\end{itemize}

Минимальный фрагмент из \texttt{CMakeLists.txt} проекта:

\begin{lstlisting}[style=cmake]
find_package(OpenCV 4.5 REQUIRED COMPONENTS core imgproc imgcodecs highgui dnn video)
find_package(Threads REQUIRED)

add_library(videotest_core STATIC
    src/logger.cpp
    src/pipeline.cpp
    src/track/tracker.cpp
    ...)
target_include_directories(videotest_core PUBLIC src "${ONNXRUNTIME_ROOT}/include")
target_link_libraries(videotest_core PUBLIC
    ${OpenCV_LIBS} Threads::Threads ${ONNXRUNTIME_LIBRARY})

add_executable(videotest src/main.cpp)
target_link_libraries(videotest PRIVATE videotest_core)
\end{lstlisting}

\subsection{Флаги компилятора, которые стоит понимать}

\begin{center}
\begin{tabular}{ll}
\toprule
Флаг & Смысл \\
\midrule
\texttt{-std=c++17} & версия стандарта языка \\
\texttt{-O2}, \texttt{-O3} & уровень оптимизации \\
\texttt{-g} & отладочная информация (для \texttt{gdb}) \\
\texttt{-Wall -Wextra} & включить большинство предупреждений \\
\texttt{-fsanitize=address} & AddressSanitizer: ловит выход за границы \\
\texttt{-fsanitize=undefined} & UBSan: ловит неопределённое поведение \\
\texttt{-Werror} & превратить предупреждения в ошибки \\
\bottomrule
\end{tabular}
\end{center}

\begin{note}
\texttt{-fsanitize=address} --- самый быстрый способ найти причину
\texttt{Segmentation fault}. В этом проекте именно так была найдена ошибка
индексации в трекере: \texttt{gdb} в среде не было, но ASan указал точную
строку (\texttt{stl\_bvector.h:96}, то есть обращение к \texttt{std::vector<bool>}
за границей).
\end{note}

\subsection{Тестирование через CTest}

\texttt{enable\_testing()} в CMake и \texttt{add\_test(NAME ... COMMAND ...)}
регистрируют тесты; запуск --- \texttt{ctest --output-on-failure}. Проект
сознательно обходится без фреймворков: тесты --- это обычные программы,
возвращающие ненулевой код при ошибке.

\begin{lstlisting}[escapeinside={(*@}{@)}]
(*@\#@*)define CHECK(cond)                                                       (*@\textbackslash@*)
    do {                                                                  (*@\textbackslash@*)
        if (!(cond)) {                                                    (*@\textbackslash@*)
            std::printf("  FAIL %s:%d %s(*@\textbackslash@*)n", __FILE__, __LINE__, (*@\#@*)cond);  (*@\textbackslash@*)
            ++g_failed;                                                   (*@\textbackslash@*)
        }                                                                 (*@\textbackslash@*)
    } while (0)
\end{lstlisting}

Макрос \texttt{\#cond} --- приём препроцессора: аргумент превращается в
строку, поэтому в отчёте видно само условие, а не только номер строки.

\subsubsection{Главное правило: тест должен падать на ошибке}

Бесполезный тест --- тот, который проходит всегда. Поэтому после написания
теста его проверяют \emph{обратным} экспериментом: временно возвращают
ошибочный код и убеждаются, что тест краснеет.

В проекте так проверены оба новых теста:

\begin{center}
\begin{tabular}{p{0.32\linewidth}p{0.28\linewidth}p{0.30\linewidth)}
\toprule
Что тестируется & Как ломали & Результат \\
\midrule
Сохранение \texttt{track\_id} при движении & убрали
\texttt{kalman.predict()} & 3 проваленных проверки \\
Баланс очереди & при переполнении выбрасывали
\emph{новый} элемент вместо старого & 2 проваленных проверки \\
\bottomrule
\end{tabular}
\end{center}

% ============================================================================
\section{Язык C++, который используется в проекте}

\subsection{Управление памятью: стек, куча, времена жизни}

В C++ нет сборщика мусора. Объект, созданный в автоматической памяти (на
стеке), уничтожается при выходе из области видимости; объект в динамической
памяти (в куче, \texttt{new}) живёт до явного \texttt{delete}.

\begin{lstlisting}
void f() {
    std::string s = "локальная";       // стек: живёт до конца функции
    auto* p = new std::string("куча"); // куча: живёт до delete p
    delete p;                          // иначе утечка
}   // s уничтожен автоматически
\end{lstlisting}

Ручное \texttt{new}/\texttt{delete} в современном C++ почти не применяется:
его заменяют \emph{умные указатели} и принцип RAII.

\subsection{RAII: ресурс --- это объект}

\begin{definition}
\textbf{RAII} (Resource Acquisition Is Initialization) --- приём, при котором
владение ресурсом (память, файл, блокировка, сокет) привязано к времени жизни
объекта: ресурс захватывается в конструкторе и освобождается в деструкторе.
\end{definition}

Пример из проекта --- деструктор источника видео:

\begin{lstlisting}
V4l2Source::~V4l2Source() {
    close();          // STREAMOFF -> munmap -> ::close
}
\end{lstlisting}

Порядок здесь не случаен: нельзя размечать (mmap) память уже закрытого
дескриптора. RAII делает этот порядок единственно возможным.

\subsection{Умные указатели}

\begin{center}
\begin{tabular}{p{0.22\linewidth}p{0.34\linewidth}p{0.34\linewidth}}
\toprule
Тип & Владение & Когда применяется \\
\midrule
\texttt{std::unique\_ptr<T>} & единоличное & ресурс принадлежит одному объекту \\
\texttt{std::shared\_ptr<T>} & разделяемое, со счётчиком ссылок &
данные нужны нескольким владельцам \\
\texttt{std::weak\_ptr<T>} & не владеет & наблюдение без продления жизни \\
\bottomrule
\end{tabular}
\end{center}

\subsubsection{Почему в конвейере именно \texttt{shared\_ptr}}

Кадр нужен одновременно в двух местах: он лежит в очереди, и на него же
ссылается результат обработки. При этом заранее неизвестно, кто отпустит его
последним.

\begin{lstlisting}
struct Frame {
    cv::Mat image;
    std::uint64_t id = 0;
    std::chrono::steady_clock::time_point captured;
};
struct Result {
    std::shared_ptr<const Frame> frame;   // разделяемое владение, только чтение
    std::vector<Instance> instances;
};

using FramePtr  = std::shared_ptr<Frame>;
using ResultPtr = std::shared_ptr<Result>;
\end{lstlisting}

Копия \texttt{shared\_ptr} увеличивает счётчик ссылок атомарно; уничтожение
последней копии освобождает объект. Пиксели при этом \emph{не копируются}:
\texttt{cv::Mat} внутри \texttt{Frame} сам является дешёвым дескриптором с
разделяемым буфером.

\subsubsection{Почему \texttt{shared\_ptr<const Frame>}}

Это ключевой приём всей потокобезопасности проекта. Тип \texttt{const Frame}
делает кадр \emph{неизменяемым} для всех, кто держит указатель. Пока
отрисовывающий поток читает пиксели, обрабатывающий поток уже занят следующим
кадром --- но записать в этот кадр он не может, потому что \texttt{const}.
Гонки данных нет \emph{без единого мьютекса} вокруг самих данных.

\begin{note}
\texttt{const} в C++ защищает от записи через данный указатель, но не
гарантирует неизменность объекта вообще. Здесь этого достаточно, потому что
никто больше не имеет неконстантного указателя на тот же \texttt{Frame}.
\end{note}

\subsection{Ссылки, копирование, перемещение}

\begin{itemize}
  \item \texttt{const T\&} --- ссылка без копирования, без права изменять.
        Основной способ передавать большие объекты.
  \item \texttt{T\&} --- ссылка с правом изменять (выходной параметр).
  \item \texttt{std::move(x)} --- не перемещает сам по себе, а \emph{разрешает}
        вызвать перемещающий конструктор: содержимое объекта переходит в новый,
        исходный остаётся в допустимом, но неопределённом состоянии.
\end{itemize}

\begin{lstlisting}
// возврат через выходной параметр: без копирования и без исключений
bool grab(FramePtr& out, std::string& error);
\end{lstlisting}

Такой стиль (функция возвращает \texttt{bool}, результат --- через ссылку,
текст ошибки --- тоже через ссылку) выбран потому, что он не требует
исключений на границе потоков и хорошо читается в циклах.

\subsection{Классы, виртуальные функции, интерфейсы}

Абстрактный класс \texttt{Source} задаёт контракт:

\begin{lstlisting}
class Source {
public:
    virtual ~Source() = default;
    virtual bool open(std::string& error) = 0;
    virtual bool grab(FramePtr& out, std::string& error) = 0;
    virtual int width() const = 0;
    virtual int height() const = 0;
};
\end{lstlisting}

\begin{itemize}
  \item \texttt{= 0} делает функцию \emph{чисто виртуальной}: класс нельзя
        создать, только унаследовать от него.
  \item \texttt{virtual \~{}Source() = default} обязателен: удаление объекта
        производного класса через указатель на базовый без виртуального
        деструктора --- неопределённое поведение.
  \item Вызов виртуальной функции идёт через таблицу виртуальных функций
        (vtable), то есть стоит одного косвенного перехода. Для вызова раз в
        кадр это несущественно.
\end{itemize}

\subsection{Шаблоны}

Очередь --- шаблон класса, потому что одна и та же логика нужна для
\texttt{FramePtr} и \texttt{ResultPtr}:

\begin{lstlisting}
template <class T>
class Queue {
public:
    std::size_t push(T value);
    template <class Rep, class Period>
    std::optional<T> pop(std::chrono::duration<Rep, Period> timeout);
};
\end{lstlisting}

Второй шаблон (шаблонный метод) нужен, чтобы принимать любой тип длительности:
\texttt{100ms}, \texttt{std::chrono::milliseconds(10)} и т.п.

\subsection{\texttt{std::optional}: отсутствие значения как тип}

\begin{lstlisting}
std::optional<T> pop(...);   // либо значение, либо «данных нет»
\end{lstlisting}

Возвращать <<магическое>> значение (например, пустой указатель) можно, но
тип не говорит об этом вызывающему. \texttt{std::optional} делает отсутствие
значения частью типа, и компилятор заставляет его обработать.

\subsection{Идиома pimpl, или сокрытие зависимости}

Заголовок \texttt{segmenter.hpp} не включает \texttt{onnxruntime\_cxx\_api.h}:
этот заголовок тяжёлый и содержит макросы, которые не должны протекать во все
единицы трансляции. Вместо этого объявлена неполная структура:

\begin{lstlisting}
struct OrtHolder;          // объявление без определения
OrtHolder* ort_ = nullptr; // указатель на неё
\end{lstlisting}

Определение живёт в \texttt{.cpp}:

\begin{lstlisting}
struct Segmenter::OrtHolder {
    Ort::Env env{ORT_LOGGING_LEVEL_ERROR, "videotest"};
    Ort::SessionOptions options;
    std::unique_ptr<Ort::Session> session;
    Ort::MemoryInfo memory{"Cpu", OrtArenaAllocator, 0, OrtMemTypeDefault};
};
\end{lstlisting}

\begin{note}
Здесь ручное \texttt{new}/\texttt{delete} (в деструкторе \texttt{Segmenter}).
Это осознанный компромисс. С \texttt{std::unique\_ptr<OrtHolder>} требуется
полное определение типа там, где генерируется код уничтожения, то есть в
\texttt{.cpp}. Технически это решается объявлением деструктора в заголовке и
его определением в \texttt{.cpp} (что и сделано), но тогда
\texttt{unique\_ptr} не даёт выигрыша в читаемости: деструктор всё равно
написан вручную. Вариант с указателем явнее показывает, что владение
единоличное и освобождается в одном месте.
\end{note}

\subsection{Обработка ошибок и границы потоков}

В C++ исключение, покинувшее поток, вызывает \texttt{std::terminate()}.
Поэтому на границе рабочего потока исключения \emph{переводятся в код
возврата}:

\begin{lstlisting}
try {
    ... session->Run(...);
} catch (const Ort::Exception& e) {
    error = std::string("ONNX Runtime: ") + e.what();
    return false;      // исключение не уходит в std::thread
}
\end{lstlisting}

% ============================================================================
\section{Данные, ходящие между потоками}

\subsection{Три структуры и их инварианты}

\begin{center}
\begin{tabular}{p{0.16\linewidth}p{0.76\linewidth}}
\toprule
Тип & Инварианты \\
\midrule
\texttt{Frame} & \texttt{image} --- BGR, \texttt{CV\_8UC3}; \texttt{captured} ---
момент получения от драйвера, монотонные часы. \\
\texttt{Instance} & \texttt{box} в целых пикселях исходного кадра;
\texttt{mask} либо пустая, либо \texttt{CV\_8UC1} размера кадра со значениями
0/255; \texttt{track\_id = -1} до первого присвоения. \\
\texttt{Result} & \texttt{frame} --- неизменяемый указатель на кадр; вектор
\texttt{instances} отсортирован по убыванию уверенности. \\
\bottomrule
\end{tabular}
\end{center}

\begin{note}
Почему \texttt{steady\_clock}, а не \texttt{system\_clock}: последний может быть
переведён назад (синхронизация времени, переход на летнее время), и тогда
<<задержка>> окажется отрицательной. \texttt{steady\_clock} гарантированно
монотонен.
\end{note}

\begin{remark}
Почему координаты целые: рисование всё равно округляет, а целочисленные
операции дешевле. \texttt{cv::Rect2f} применяется только внутри
преобразований, где важна дробная точность.
\end{remark}

% ============================================================================
\section{Очередь с вытеснением: механизм и обоснование}

\subsection{Постановка задачи}

Нужен буфер между производителем (камера, 30 кадров/с, ждать нельзя) и
потребителем (модель, время обработки переменное). Требования:

\begin{enumerate}
  \item производитель никогда не блокируется;
  \item потребитель получает \emph{свежие} данные;
  \item задержка <<попадания в обработку>> ограничена константой;
  \item остановка программы не должна зависать на ожидании данных.
\end{enumerate}

Классическая очередь FIFO без ограничения размера нарушает пункты 2 и 3:
если модель медленнее камеры, длина очереди растёт, и обрабатывается кадр
десятисекундной давности.

\subsection{Политика: вытеснение самого старого}

\begin{definition}[Очередь с вытеснением]
Очередь ёмкости $C$: операция \textsc{push} всегда завершается успешно. Если
очередь заполнена, предварительно удаляется элемент, поступивший раньше всех.
Операция \textsc{pop} возвращает самый старый элемент либо, при таймауте, ничего.
\end{definition}

Точное соотношение баланса: в любой момент
\[
  N \;=\; M \;+\; D \;+\; |Q|,
\]
где $N$ --- сколько раз вызван \textsc{push}, $M$ --- сколько раз успешно
вызван \textsc{pop}, $D$ --- сколько элементов вытеснено (это число возвращает
\textsc{push}), $|Q|$ --- текущее содержимое очереди. После того как
производитель закончил и очередь опустела, $|Q| = 0$ и
\[
  D = N - M .
\]
Именно это тождество проверяет тест \texttt{queue\_tests}.

\subsection{Реализация}

\begin{lstlisting}
template <class T>
class Queue {
public:
    explicit Queue(std::size_t capacity)
        : capacity_(capacity == 0 ? 1 : capacity) {}

    std::size_t push(T value) {
        std::size_t dropped = 0;
        {
            std::lock_guard lock(mutex_);
            if (items_.size() >= capacity_) {
                items_.pop_front();
                ++dropped;
            }
            items_.push_back(std::move(value));
        }
        not_empty_.notify_one();
        return dropped;
    }

    template <class Rep, class Period>
    std::optional<T> pop(std::chrono::duration<Rep, Period> timeout) {
        std::unique_lock lock(mutex_);
        not_empty_.wait_for(lock, timeout, [this] { return !items_.empty(); });
        if (items_.empty()) return std::nullopt;
        T value = std::move(items_.front());
        items_.pop_front();
        return value;
    }
private:
    mutable std::mutex mutex_;
    std::condition_variable not_empty_;
    std::deque<T> items_;
    const std::size_t capacity_;
};
\end{lstlisting}

Разберём каждое решение.

\subsubsection{Почему \texttt{std::deque}, а не \texttt{std::vector}}

Вытеснение удаляет элемент из \emph{начала}, добавление идёт в \emph{конец}.
\texttt{std::deque} поддерживает обе операции за $O(1)$; \texttt{std::vector}
удаление из начала делает за $O(n)$. При ёмкости 2--4 это неважно по
абсолютному времени, но выбор структуры должен соответствовать операциям.

\subsubsection{Почему \texttt{std::mutex} и \texttt{condition\_variable}}

Мьютекс защищает инвариант <<\texttt{items\_} согласован>>. Условная переменная
позволяет потребителю \emph{спать}, а не крутиться в цикле, расходуя ядро.

\subsubsection{Почему \texttt{unique\_lock}, а не \texttt{lock\_guard}}

\texttt{wait\_for} должен \emph{отпустить} мьютекс на время сна и снова взять
его после пробуждения. \texttt{std::unique\_lock} это умеет,
\texttt{std::lock\_guard} --- нет.

\subsubsection{Почему предикат обязателен}

\begin{lstlisting}
not_empty_.wait_for(lock, timeout, [this] { return !items_.empty(); });
\end{lstlisting}

Условная переменная может проснуться без уведомления (\emph{ложное
пробуждение}) и может быть разбужена не тем событием. Без предиката код
продолжил бы работу с пустой очередью. С предикатом \texttt{wait\_for}
возвращается только когда условие истинно либо истёк таймаут.

\subsubsection{Почему уведомление вне блокировки}

\begin{lstlisting}
{
    std::lock_guard lock(mutex_);
    ...
}                 // блокировка снята
not_empty_.notify_one();   // уведомляем после
\end{lstlisting}

Если уведомить под мьютексом, разбуженный поток сразу же упрётся в
заблокированный мьютекс (\emph{thundering herd} / лишнее переключение).
Уведомление после снятия блокировки этого избегает.

\subsubsection{Почему таймаут в \texttt{pop}}

Без таймаута поток-обработчик, ожидая кадр, никогда не узнал бы о том, что
программу просят остановить. Таймаут 100 мс означает: <<даже если данных нет,
проснись и проверь флаг остановки>>.

\subsection{Ёмкость 2: почему именно столько}

\begin{itemize}
  \item Ёмкость 1: пока потребитель обрабатывает кадр, производитель уже
        положил следующий --- один кадр гарантированно ждёт. Любая задержка
        сразу приводит к вытеснению.
  \item Ёмкость 2: допускает кратковременную неравномерность (один кадр в
        работе, один в запасе), при этом задержка ограничена сверху временем
        двух кадров.
  \item Ёмкость $\gg 2$: задержка растёт пропорционально, смысл вытеснения
        теряется.
\end{itemize}

\begin{note}
Число вытеснений выведено в HUD как \texttt{dropped}. Это прямой индикатор:
если счётчик растёт, модель не успевает за камерой, и на экране видно меньше
кадров, чем снимает камера. Именно так и должно быть --- обработка не
тормозит приём.
\end{note}

% ============================================================================
\section{Потоки и синхронизация}

\subsection{Схема}

\begin{center}
\begin{tabular}{lll}
\toprule
Поток & Что делает & Куда пишет \\
\midrule
1. Захват & \texttt{source.grab()} в цикле & \texttt{Queue<FramePtr>}(2) \\
2. Обработка & \texttt{pop()} $\to$ \texttt{infer()} $\to$ \texttt{tracker.update()} &
\texttt{Queue<ResultPtr>}(2) \\
Главный & \texttt{pop(10\,ms)} $\to$ отрисовка $\to$ \texttt{imshow} & экран \\
\bottomrule
\end{tabular}
\end{center}

\subsection{Порядок остановки --- источник тонкой ошибки}

\begin{lstlisting}
void Pipeline::stop() {
    running_.store(false, std::memory_order_relaxed);  // 1. снять флаг
    if (capture_thread_.joinable()) capture_thread_.join();   // 2. дождаться
    if (process_thread_.joinable()) process_thread_.join();
}
\end{lstlisting}

Если поменять местами, получится \emph{дедлок}: \texttt{join} ждёт завершения
потока, а поток крутится по флагу, который ещё не снят. В первой версии эта
ошибка присутствовала и проявлялась как \texttt{terminate called without an
active exception}.

\begin{note}
Почему \texttt{terminate}: \texttt{std::thread}, который остаётся
\emph{joinable} при разрушении, вызывает \texttt{std::terminate}. Разрушение
происходило по выходу из \texttt{main} при уже завершившихся циклах.
\end{note}

\subsection{Атомарные счётчики}

Метрики (получено, обработано, вытеснено) читаются главным потоком и
пишутся рабочими. Здесь достаточно \texttt{std::atomic<std::uint64\_t>}:
операция \texttt{fetch\_add} не требует блокировки. Это дешевле мьютекса
и не создаёт конкуренции на кадр.

\subsection{Остановка по Ctrl+C}

Обработчик сигнала обязан быть \emph{асинхронно-безопасным}: в нём нельзя
брать мьютексы, вызывать \texttt{printf}, выделять память. Единственное
допустимое действие --- записать в атомарный флаг:

\begin{lstlisting}
extern "C" void on_signal(int) {
    stop_flag().store(true, std::memory_order_relaxed);
}
\end{lstlisting}

Рабочие циклы сами проверяют флаг: \texttt{poll()} возвращается по таймауту
500 мс, \texttt{Queue::pop()} --- по таймауту 100 мс.

% ============================================================================
\section{Захват видео через V4L2}

\subsection{Что такое V4L2}

Video4Linux2 --- набор системных вызовов \texttt{ioctl} для работы с
видеоустройствами в Linux. Камера представляется файлом \texttt{/dev/videoX};
управление идёт командами \texttt{VIDIOC\_*}.

\subsection{Последовательность вызовов}

\begin{lstlisting}[style=shell]
open(/dev/videoX, O_RDWR | O_NONBLOCK)
  VIDIOC_QUERYCAP       есть ли захват и потоковая передача
  VIDIOC_S_FMT          задать формат и размер
  VIDIOC_S_PARM         запросить частоту кадров
  VIDIOC_REQBUFS        выделить N буферов драйвера (MMAP)
  VIDIOC_QUERYBUF + mmap  спроецировать буферы в своё адресное пространство
  VIDIOC_QBUF x N       отдать буферы драйверу
  VIDIOC_STREAMON       запустить поток
цикл: poll() -> VIDIOC_DQBUF -> конвертация -> VIDIOC_QBUF
\end{lstlisting}

\subsection{Тонкости, без которых не работает}

\begin{enumerate}
  \item \textbf{\texttt{device\_caps}, а не \texttt{capabilities}.} У камеры с
        микрофоном флаг \texttt{capabilities} описывает всё устройство. Реальные
        возможности конкретного узла лежат в \texttt{device\_caps}, и проверка
        только по \texttt{capabilities} ложно отвергает часть камер.
  \item \textbf{\texttt{poll} перед \texttt{DQBUF}.} Файл открыт с
        \texttt{O\_NONBLOCK}, поэтому \texttt{DQBUF} без данных вернул бы
        \texttt{EAGAIN}, и пришлось бы крутиться в busy-loop, съедая ядро.
  \item \textbf{\texttt{QBUF} обязателен всегда.} После \texttt{DQBUF} буфер
        принадлежит программе. Не вернув его, через \texttt{buffer\_count}
        кадров драйверу нечего заполнять, и камера встанет.
  \item \textbf{\texttt{bytesperline} (шаг строки).} Длина строки в памяти может
        быть больше $\text{width} \times \text{bpp}$ из-за выравнивания.
        Матрица создаётся с явным шагом
        (\texttt{cv::Mat(height, width, CV\_8UC2, data, step)}); без него
        картинка <<поедет>> по диагонали.
  \item \textbf{\texttt{EINTR}.} Системный вызов может быть прерван сигналом
        (например, нашим SIGINT). Это не ошибка устройства, поэтому
        \texttt{ioctl} повторяется (в проекте --- шаблонная обёртка
        \texttt{xioctl}, принимающая переменное число аргументов):
        \begin{lstlisting}
template <class... Args>
int xioctl(int fd, unsigned long request, Args... args) {
    int result;
    do {
        result = ::ioctl(fd, request, args...);
    } while (result == -1 && errno == EINTR);
    return result;
}
        \end{lstlisting}
\end{enumerate}

\subsection{Почему mmap, а не read}

При \texttt{read()} драйвер копирует кадр в буфер пользователя при каждом
вызове. При mmap память драйвера отображается в адресное пространство
процесса: копия остаётся только при конвертации формата (её делает OpenCV).
На 1280$\times$720 и 30 кадрах в секунду это заметная разница.

\subsection{MJPG и YUYV}

\begin{center}
\begin{tabular}{p{0.12\linewidth}p{0.38\linewidth}p{0.40\linewidth}}
\toprule
Формат & Хранение & Следствия \\
\midrule
MJPG & каждый кадр --- JPEG & меньше по шине USB, но требует декодирования
(единицы мс); нет потерь в \texttt{Mat}, потери в JPEG \\
YUYV & 2 байта на пиксель, подвыборка цветности & не требует декодирования,
но по USB~2.0 не всегда влезает 1280$\times$720 при 30 кадрах/с \\
\bottomrule
\end{tabular}
\end{center}

\begin{remark}
Формат в проекте задаётся \emph{явно} (\texttt{--format MJPG|YUYV}).
Автоподбор убран сознательно: он скрывает, что реально поддерживает устройство.
Если драйвер подменил формат, программа выдаёт ошибку с подсказкой
\texttt{v4l2-ctl --list-formats-ext}.
\end{remark}

% ============================================================================
\section{Модель: от тензора к объектам}

\subsection{Препроцессинг: letterbox}

Сеть ждёт квадратный вход $S \times S$ (обычно $S = 640$), значения
$\mathbb{R}^{3 \times S \times S}$ в диапазоне $[0,1]$, порядок каналов RGB.
Кадр имеет размер $W \times H$.

\subsubsection{Математика}

Вычисляется масштаб
\[
  s = \min\left(\frac{S}{W}, \frac{S}{H}\right),
\]
новый размер $W' = \lfloor sW \rceil$, $H' = \lfloor sH \rceil$, и отступы
\[
  p_x = \frac{S - W'}{2}, \qquad p_y = \frac{S - H'}{2}.
\]

\begin{definition}[Прямое преобразование]
Пиксель кадра $(u,v)$ отображается в точку сети
\[
  x = s u + p_x, \qquad y = s v + p_y.
\]
\end{definition}

\begin{definition}[Обратное преобразование]
По координатам сети $(x,y)$ координата кадра восстанавливается как
\[
  u = \frac{x - p_x}{s}, \qquad v = \frac{y - p_y}{s}.
\]
\end{definition}

Пара $(s, p_x, p_y)$ --- единственное, что нужно сохранить при прямом ходе,
чтобы правильно вернуть координаты обратно. Именно поэтому
\texttt{preprocess()} возвращает \texttt{scale} и \texttt{pad} по ссылке.

\begin{note}
Почему нельзя просто \texttt{resize} в квадрат: при разных $W$ и $H$ это
даёт \emph{разные} масштабы по осям, и обратное преобразование $x$ и $y$
должно выполняться с разными коэффициентами. Забыть об одном из них ---
классическая ошибка, дающая сдвиг, растущий к краю кадра.
\end{note}

Отступы заполняются значением 114: это стандартное значение YOLO при
обучении, и сеть к нему привыкла.

\subsubsection{Приведение к входу сети}

\begin{equation}
  \text{blob}[c][y][x] = \frac{\text{canvas}[y][x][c']}{255},
  \qquad c' = \begin{cases}
    2 - c, & \text{для RGB (swapRB)},\\
    c, & \text{иначе},
  \end{cases}
\end{equation}
где $c \in \{0,1,2\}$ --- индекс канала в NCHW. Функция
\texttt{cv::dnn::blobFromImage} выполняет и перестановку осей HWC$\to$NCHW,
и деление на 255, и перевод BGR$\to$RGB.

\subsection{Выходы модели}

\subsubsection{Формат \texttt{raw} (YOLOv8-seg, YOLO11-seg)}

\[
  \text{output0} \in \mathbb{R}^{1 \times (4 + n_c + 32) \times N},
  \qquad
  \text{output1} \in \mathbb{R}^{1 \times 32 \times h \times w},
\]
где $n_c$ --- число классов, $N$ --- число якорей (8400 при $S=640$),
$h, w$ --- размер карты прототипов (160 при $S = 640$).

Раскладка \emph{канал-мажорная}: значение канала $c$ для якоря $a$ лежит по
индексу
\[
  \text{idx}(c,a) = c \cdot N + a .
\]

Столбцы: $c_x, c_y, w, h$, затем $n_c$ вероятностей классов, затем 32
коэффициента масок.

\subsubsection{Формат \texttt{nms} (YOLO26, end-to-end)}

\[
  \text{output0} \in \mathbb{R}^{1 \times M \times (6 + 32)}, \qquad M = 300,
\]
столбцы: $x_1, y_1, x_2, y_2$, уверенность, номер класса, 32 коэффициента.
Подавление пересечений уже выполнено внутри графа.

\begin{note}
Повторно применять NMS к выходу \texttt{nms} нельзя: это удалит легитимно
перекрывающиеся объекты, которые сеть намеренно оставила.
\end{note}

\subsection{Декодирование и уверенность}

Для каждого якоря $a$ уверенность класса определяется как максимум по классам:
\[
  (\hat{s}, \hat{k}) = \bigl(\max_{c} p_{c,a},\; \arg\max_{c} p_{c,a}\bigr).
\]
Якорь отбрасывается, если $\hat{s} < \theta_{\text{score}}$ (порог 0.35 по
умолчанию) либо если ширина/высота неположительны.

\subsection{Подавление пересечений (NMS)}

\subsubsection{Метрика IoU}

\begin{definition}[IoU, Intersection over Union]
Для прямоугольников $A$ и $B$
\[
  \mathrm{IoU}(A,B) = \frac{|A \cap B|}{|A \cup B|}
  = \frac{|A \cap B|}{|A| + |B| - |A \cap B|}.
\]
\end{definition}

Для прямоугольников, выровненных по осям, площади вычисляются через
пересечение проекций:
\[
  |A \cap B| = \max(0, x_2^{\cap} - x_1^{\cap}) \cdot \max(0, y_2^{\cap} - y_1^{\cap}),
\]
\[
  x_1^{\cap} = \max(x_1^A, x_1^B),\quad x_2^{\cap} = \min(x_2^A, x_2^B),
\]
и аналогично по $y$. Знаменатель нулевым быть не может (площади
положительны), но в коде предусмотрена защита от деления на ноль для
вырожденных боксов.

\subsubsection{Алгоритм}

\begin{enumerate}
  \item Отсортировать детекции по убыванию уверенности.
  \item Взять лучшую непогашенную детекцию, оставить её.
  \item Погасить все следующие детекции \emph{того же класса} с
        $\mathrm{IoU} > \theta_{\text{nms}}$ (0.45 по умолчанию).
  \item Повторять, пока есть непогашенные.
\end{enumerate}

\begin{note}
Подавление \emph{class-aware} принципиально: человек и чашка часто
перекрываются. Подавление по всем классам склеило бы их в один объект.
Сложность алгоритма --- $O(n^2)$ по числу детекций после порога уверенности;
на практике $n$ невелико (десятки).
\end{note}

\subsection{Маски: линейная комбинация прототипов}

Сеть не выдаёт маску на объект. Она выдаёт:
\begin{itemize}
  \item базис из 32 <<мягких>> масок (прототипов) $P_k \in
        \mathbb{R}^{h \times w}$, $k = 1..32$;
  \item для каждого объекта --- 32 коэффициента $e_k$.
\end{itemize}

Маска объекта определяется как
\[
  M(x,y) = \sigma\Bigl(\sum_{k=1}^{32} e_k \, P_k(y,x)\Bigr),
  \qquad \sigma(z) = \frac{1}{1 + e^{-z}} .
\]

\subsubsection{Порог без вычисления сигмоиды}

Условие $\sigma(z) > t$ эквивалентно
\[
  z > \ln\frac{t}{1-t} =: \operatorname{logit}(t).
\]
При $t = 0.5$ порог равен нулю. Это позволяет заменить $h \cdot w$
вычислений экспоненты одним сравнением:

\begin{lstlisting}
const float logit = std::log(config_.mask_threshold / (1.0f - config_.mask_threshold));
cv::compare(resized, logit, binary, cv::CMP_GT);
\end{lstlisting}

\subsubsection{Приведение маски к кадру}

Прототипы имеют размер $h \times w$ (160), что соответствует входу сети
$S \times S$. Бокс объекта задан в координатах сети. Поэтому:

\begin{enumerate}
  \item линейная комбинация даёт маску $h \times w$;
  \item обрезаем её по боксу в масштабе прототипов: коэффициенты
        $\rho_x = w/S$, $\rho_y = h/S$;
  \item масштабируем обрезанное к размеру бокса \emph{в пикселях кадра};
  \item сравниваем с логит-порогом;
  \item вставляем в пустое изображение размера кадра.
\end{enumerate}

% ============================================================================
\section{Трекинг: фильтр Калмана и сопоставление}

\subsection{Зачем нужен трекер}

Детектор работает независимо на каждом кадре:
\begin{itemize}
  \item рамка <<дрожит>> из-за шума сети;
  \item на части кадров объект не детектируется вовсе (перекрытие, размытие);
  \item нет понятия <<тот же объект>> --- только множество рамок.
\end{itemize}

Трекер решает три задачи: сглаживание, предсказание при пропуске детекции и
присвоение устойчивого идентификатора.

\subsection{Модель состояния}

\begin{definition}
Вектор состояния трека
\[
  \mathbf{x} = (c_x, c_y, w, h, v_x, v_y, v_w, v_h)^\top \in \mathbb{R}^8,
\]
где первые четыре компоненты --- центр и размеры рамки, последние четыре ---
их скорости.
\end{definition}

\begin{definition}[Модель равномерного движения]
Матрица перехода
\[
  F = \begin{pmatrix}
    I_4 & I_4 \\
    0   & I_4
  \end{pmatrix} \in \mathbb{R}^{8\times 8},
\]
то есть $\mathbf{x}_{t+1} = F\mathbf{x}_t$ даёт
$c_x \leftarrow c_x + v_x$ и т.д.
\end{definition}

\begin{definition}[Модель измерений]
Детектор сообщает только положение и размер, поэтому
\[
  H = \begin{pmatrix} I_4 & 0 \end{pmatrix} \in \mathbb{R}^{4\times 8},
  \qquad \mathbf{z}_t = H\mathbf{x}_t + \mathbf{v}_t .
\]
\end{definition}

\subsection{Уравнения фильтра}

Фильтр Калмана поддерживает оценку $\hat{\mathbf{x}}$ и её ковариацию $P$.

\paragraph{Предсказание.}
\begin{align}
  \hat{\mathbf{x}}^-_t &= F \hat{\mathbf{x}}_{t-1}, \\
  P^-_t &= F P_{t-1} F^\top + Q .
\end{align}

\paragraph{Коррекция по измерению $\mathbf{z}_t$.}
\begin{align}
  K_t &= P^-_t H^\top \bigl(H P^-_t H^\top + R\bigr)^{-1}, \\
  \hat{\mathbf{x}}_t &= \hat{\mathbf{x}}^-_t + K_t\bigl(\mathbf{z}_t - H\hat{\mathbf{x}}^-_t\bigr), \\
  P_t &= (I - K_t H) P^-_t .
\end{align}

Здесь $Q$ --- ковариация шума процесса, $R$ --- ковариация шума измерений.
В проекте обе заданы диагональными:
\[
  Q = 10^{-2} I_8, \qquad R = 10^{-1} I_4, \qquad P_0 = I_8 .
\]

\subsection{Смысл матрицы усиления}

Матрица $K_t$ определяет, насколько доверять новому измерению. Если
$P^- \gg R$ (неопределённость велика, измерения точны), то $K \to I$ и оценка
почти равна измерению. Если $P^- \ll R$ (оценка уверенна, измерения шумны),
$K \to 0$ и оценка почти не меняется.

\paragraph{Почему нельзя заменить предсказание на $F\mathbf{x}$.}
Вторая строка, $P^- = FPF^\top + Q$, обязательна. Она \emph{увеличивает}
неопределённость перед новым измерением. Без неё $P$ не растёт, $K$ остаётся
крошечным, и коррекция почти не сдвигает состояние: трек <<залипает>> на
начальном положении. В проекте эта ошибка присутствовала и проявлялась так:
объект уезжал, IoU с его рамкой падал ниже порога, трек терялся, и
идентификатор пересоздавался почти каждый кадр.

\begin{note}
Практический вывод, зафиксированный тестом \texttt{tracker\_tests}: объект,
смещающийся на 15 пикселей за кадр, обязан сохранить свой идентификатор на
протяжении 25 кадров. Без вызова \texttt{cv::KalmanFilter::predict()} тест
падает.
\end{note}

\subsection{Сопоставление детекций и треков}

\begin{enumerate}
  \item Для каждой пары (детекция, трек) вычисляется $\mathrm{IoU}$.
  \item Пары с $\mathrm{IoU} \ge 0.3$ сортируются по убыванию.
  \item Жадно выбираются пары без повторов: одна детекция --- один трек.
  \item Сопоставленные треки корректируются измерением, получают
        прежний идентификатор.
  \item Несопоставленные детекции создают новые треки с новым идентификатором.
  \item Треки без детекции увеличивают счётчик пропусков; при превышении
        \texttt{max\_age} = 15 удаляются.
\end{enumerate}

\begin{note}
Жадное сопоставление --- упрощение по сравнению с венгерским алгоритмом
(оптимальным по суммарному IoU). На 10--50 объектах разница в качестве мала,
а код проще на порядок. Замена локализована в одной функции.
\end{note}

\subsection{Обоснование выбора порогов}

\begin{itemize}
  \item $\mathrm{IoU} \ge 0.3$: рамки соседних кадров при 30 кадрах/с
        перекрываются сильнее; порог 0.3 отсекает <<соседние>> объекты, но
        терпит дрожание.
  \item \texttt{max\_age} = 15: полсекунды при 30 кадрах/с. Достаточно, чтобы
        пережить кратковременное перекрытие, но не настолько, чтобы <<держать>>
        ушедший объект.
\end{itemize}

% ============================================================================
\section{Метрики качества}

\subsection{Метрики детекции}

\subsubsection{Точность, полнота, $F_1$}

Для детекций вводится понятие совпадения: детекция засчитывается как верная
(TP), если её IoU с размеченной рамкой превышает порог (обычно 0.5) и класс
совпадает.

\begin{center}
\begin{tabular}{lll}
\toprule
Величина & Определение & Смысл \\
\midrule
Precision & $\dfrac{TP}{TP + FP}$ & доля верных среди предсказанных \\
Recall & $\dfrac{TP}{TP + FN}$ & доля найденных среди размеченных \\
$F_1$ & $\dfrac{2PR}{P+R}$ & гармоническое среднее \\
\bottomrule
\end{tabular}
\end{center}

\subsubsection{AP и mAP}

Кривая \emph{precision--recall} строится варьированием порога уверенности.
Average Precision --- площадь под ней:
\[
  \mathrm{AP} = \int_0^1 p(r)\, dr .
\]
На практике интеграл заменяется суммой по точкам интерполяции:
\[
  \mathrm{AP} \approx \sum_k (r_{k+1} - r_k)\, \tilde{p}(r_{k+1}),
\]
где $\tilde{p}(r) = \max_{r' \ge r} p(r')$ --- монотонная огибающая.

\begin{definition}[mAP]
\[
  \mathrm{mAP} = \frac{1}{n_c} \sum_{c=1}^{n_c} \mathrm{AP}_c ,
\]
среднее AP по классам. Порог IoU фиксирует вариант метрики
(\texttt{mAP@0.5}, \texttt{mAP@0.5:0.95} --- усреднение по порогам
0.5, 0.55, \dots, 0.95).
\end{definition}

\subsubsection{Метрики сегментации}

Для масок IoU считается по пикселям:
\[
  \mathrm{IoU}_{\text{mask}}(A,B)
  = \frac{|\{p : A(p) = 1 \wedge B(p) = 1\}|}
         {|\{p : A(p) = 1 \vee B(p) = 1\}|}.
\]
\begin{itemize}
  \item \textbf{Dice} (коэффициент Сёренсена):
        $\mathrm{Dice} = \dfrac{2|A \cap B|}{|A| + |B|}$,
        связано с IoU как $\mathrm{Dice} = \dfrac{2\,\mathrm{IoU}}{1 + \mathrm{IoU}}$.
  \item \textbf{mIoU} --- среднее IoU по классам (semantic segmentation).
  \item \textbf{Mask mAP} --- AP, где совпадение определяется по IoU масок.
\end{itemize}

\subsection{Метрики трекинга}

\subsubsection{MOTA}

\[
  \mathrm{MOTA} = 1 - \frac{\sum_t \bigl(FN_t + FP_t + IDSW_t\bigr)}{\sum_t GT_t},
\]
где $FN$ --- пропущенные, $FP$ --- ложные, $IDSW$ --- переключения
идентификатора, $GT$ --- число размеченных объектов. Значение может быть
отрицательным при большом числе ошибок.

\subsubsection{IDF1 и MOTP}

\[
  \mathrm{IDF1} = \frac{2\,\mathrm{IDTP}}{2\,\mathrm{IDTP} + \mathrm{IDFP} + \mathrm{IDFN}},
\]
\[
  \mathrm{MOTP} = \frac{\sum_{i \in \mathrm{TP}} \mathrm{IoU}_i}{|\mathrm{TP}|}
\]
--- среднее совпадение для верно сопоставленных объектов (мера точности
локализации).

\subsection{Метрики реального времени}

\begin{center}
\begin{tabular}{p{0.24\linewidth}p{0.68\linewidth}}
\toprule
Метрика & Определение \\
\midrule
Задержка (latency) & $t_{\text{разметка}} - t_{\text{получение кадра}}$ \\
Пропускная способность & число кадров, обработанных в секунду \\
Доля вытеснений & $\dfrac{D}{N}$ --- показывает, насколько модель отстаёт \\
\bottomrule
\end{tabular}
\end{center}

\begin{note}
Ключевое свойство правильной схемы: при перегрузке \emph{средняя задержка не
растёт}. В проекте измерялось $640\times480$ с моделью-заглушкой: спокойный
режим (30 кадров/с) --- 42 мс, перегрузка (200 кадров/с) --- 23 мс. В
перегрузке задержка \emph{меньше}, потому что конвейер не ждёт следующего
кадра камеры: он обрабатывает непрерывно, а устаревшее вытесняется.

Числа задержки включают путь <<получено от драйвера $\to$ разметка готова>>.
В спокойном режиме значительная их часть --- это ожидание, пока кадр
полежит в очереди: он приходит раз в 33 мс, а обрабатывается за 10 мс.
\end{note}

% ============================================================================
\section{Декодирование и отрисовка: практические детали}

\subsection{Почему маски рисуются одним слоем}

Если накладывать каждую маску отдельно через \texttt{addWeighted}, в местах
перекрытия альфа складывается:
\[
  C_{\text{out}} = \alpha C_1 + (1-\alpha)\bigl(\alpha C_2 + (1-\alpha) C_0\bigr)
  \neq \alpha C_1 + (1-\alpha) C_0 ,
\]
и пятно темнеет. В проекте маски красятся в отдельный слой, который
смешивается с кадром \emph{один} раз.

\subsection{Почему \texttt{findContours} на копии}

\texttt{cv::findContours} модифицирует входную матрицу (историческая
особенность OpenCV). Поскольку \texttt{Instance.mask} может читаться другим
потоком, контуры строятся по копии ROI:

\begin{lstlisting}
cv::findContours(instance.mask(roi).clone(), contours, cv::RETR_EXTERNAL,
                 cv::CHAIN_APPROX_SIMPLE);
\end{lstlisting}

\subsection{Принцип <<нейтральной разметки>>}

Если модели нет, вектор \texttt{instances} пуст, и функция \texttt{draw\_overlay}
возвращает обычную копию кадра. Правая половина склейки оказывается
побитово равна левой. Отдельного режима <<без модели>> нет --- это важно:
конвейер всегда работает одинаково, а отсутствие модели --- частный случай
отсутствия объектов.

% ============================================================================
\section{Как читать и проверять многопоточный код}

Многопоточные ошибки отличаются от обычных тем, что проявляются
\emph{недетерминированно}. Ниже --- практические приёмы, применимые к этому
проекту.

\subsection{Правила, по которым устроен код}

\begin{enumerate}
  \item \textbf{Один владелец --- один поток.} \texttt{Queue} принадлежит
        конкретным потокам и защищена изнутри. Наружу она не отдаёт ссылки на
        своё содержимое.
  \item \textbf{Данные между потоками неизменяемы.}
        \texttt{shared\_ptr<const Frame>} --- компилятор сам не даст записать.
  \item \textbf{Состояние, разделяемое на чтение, --- атомарное.} Счётчики
        \texttt{std::atomic} читаются из HUD без блокировок.
  \item \textbf{Никаких <<флагов на честном слове>>.} Флаг остановки ---
        \texttt{std::atomic<bool>}, а не \texttt{bool}.
\end{enumerate}

\subsection{Типичные ошибки и как их узнать}

\begin{center}
\begin{tabular}{p{0.24\linewidth}p{0.30\linewidth}p{0.38\linewidth}}
\toprule
Ошибка & Проявление & Как обнаружить \\
\midrule
Гонка данных (data race) & <<плавающие>> результаты, редкие падения &
ThreadSanitizer (\texttt{-fsanitize=thread}) \\
Использование после освобождения & падение в случайном месте &
AddressSanitizer \\
Дедлок & программа висит на выходе & дамп стеков, \texttt{gdb -p} \\
Ложное пробуждение без предиката & падение в \texttt{pop()} & проверка
предиката в \texttt{wait} \\
\texttt{joinable} поток при разрушении & \texttt{terminate} & порядок
флага и \texttt{join} \\
\bottomrule
\end{tabular}
\end{center}

\subsection{Как воспроизвести проблему производительности}

Приёмы, использованные в проекте:

\begin{itemize}
  \item \textbf{Счётчик вытеснений.} Если \texttt{dropped} растёт, обработка
        не успевает. Это не ошибка, а диагностика.
  \item \textbf{Искусственная перегрузка.} Флаг
        \texttt{--fps 200} у синтетического источника создаёт заведомо
        избыточный поток кадров: на нём проверяется, что задержка не растёт.
  \item \textbf{Явные метрики в конце прогона.} Первая, средняя и
        максимальная задержка --- три числа, которых достаточно, чтобы
        отличить рабочую схему от деградирующей.
\end{itemize}

\begin{note}
Проверять надо не <<средний FPS>>, а поведение при перегрузке. Система,
которая в среднем быстра, но копит очередь, непригодна для реального
времени: задержка в ней растёт линейно и рано или поздно становится
недопустимой.
\end{note}

\section{Библиотеки и фреймворки}

\subsection{OpenCV}

Назначение: обработка изображений и видео. Используемые модули:

\begin{center}
\begin{tabular}{ll}
\toprule
Модуль & Что даёт проекту \\
\midrule
\texttt{core} & \texttt{cv::Mat}, базовые типы, арифметика \\
\texttt{imgproc} & \texttt{resize}, \texttt{findContours}, \texttt{drawContours},
                 \texttt{addWeighted} \\
\texttt{imgcodecs} & \texttt{imdecode} (JPEG из камеры), \texttt{imwrite} \\
\texttt{highgui} & \texttt{imshow}, \texttt{waitKey} \\
\texttt{dnn} & \texttt{blobFromImage} \\
\texttt{video} & \texttt{cv::KalmanFilter} \\
\bottomrule
\end{tabular}
\end{center}

\begin{note}
Модуль \texttt{videoio} не используется: видео берётся напрямую из V4L2.
Это осознанный выбор --- \texttt{VideoCapture} скрывает формат и negotiated
параметры, а цель проекта --- увидеть механику.
\end{note}

\subsection{ONNX Runtime}

Назначение: исполнение моделей в формате ONNX на CPU. Основные понятия:

\begin{itemize}
  \item \texttt{Ort::Env} --- окружение (логирование, потоки), одно на процесс;
  \item \texttt{Ort::Session} --- загруженная модель;
  \item \texttt{Ort::Value} --- тензор; создаётся на уже существующей памяти,
        поэтому данные не копируются;
  \item \texttt{SessionOptions} --- уровень оптимизации графа, число потоков.
\end{itemize}

\begin{note}
Почему \texttt{CreateTensor} на памяти \texttt{cv::Mat}: входной blob уже
лежит непрерывным массивом \texttt{float}. Копировать его в тензор --- лишние
мегабайты на каждый кадр.
\end{note}

\subsection{Прочие зависимости}

\begin{center}
\begin{tabular}{ll}
\toprule
Библиотека & Роль \\
\midrule
\texttt{pthreads} (\texttt{Threads::Threads}) & \texttt{std::thread}, мьютексы \\
V4L2 (заголовки ядра) & \texttt{linux/videodev2.h}, \texttt{sys/ioctl.h} \\
\bottomrule
\end{tabular}
\end{center}

% ============================================================================
\section{Инструменты разработки и отладки}

\subsection{Что искать при падениях}

\begin{center}
\begin{tabular}{p{0.30\linewidth}p{0.64\linewidth}}
\toprule
Симптом & Что это обычно значит \\
\midrule
\texttt{Segmentation fault} & разыменование пустого/висячего указателя,
                              выход за границы массива \\
\texttt{terminate called without an active exception} & \texttt{joinable}
        \texttt{std::thread} при разрушении, либо исключение в потоке \\
Зависание на выходе & \texttt{join} без снятия флага остановки \\
Рост памяти & утечка \texttt{new} без \texttt{delete}, цикл ссылок
              \texttt{shared\_ptr} \\
<<Пустая разметка>> & неверный layout выходов, порог уверенности,
        несовпадение числа классов \\
\bottomrule
\end{tabular}
\end{center}

\subsection{Инструменты}

\begin{itemize}
  \item \textbf{AddressSanitizer}: \texttt{-fsanitize=address -g}. Указывает
        строку и тип ошибки памяти. В этом проекте нашёл выход за границы
        \texttt{std::vector<bool>}.
  \item \textbf{UBSan}: \texttt{-fsanitize=undefined}. Ловит переполнения
        знаковых, сдвиги, выравнивание.
  \item \textbf{ThreadSanitizer}: \texttt{-fsanitize=thread}. Ищет гонки
        данных; полезен при изменении очереди.
  \item \textbf{valgrind}: медленнее, но умеет профилировать память.
  \item \textbf{perf}: профилирование CPU.
\end{itemize}

\subsection{Проверка без камеры и без весов}

\begin{lstlisting}[style=shell]
./scripts/smoke_test.sh
\end{lstlisting}

Скрипт выполняет пять шагов: сборка, тесты, прогон с моделью-заглушкой,
проверка контракта <<пустые маски $\Rightarrow$ половины совпадают>>,
корректная остановка по SIGINT. Модель-заглушка генерируется скриптом
\texttt{make\_stub\_model.py}: она имеет те же формы тензоров, что и настоящая,
но случайные веса. Этого достаточно для проверки декодера, масок, трекера и
отрисовки.

% ============================================================================
\section{Разбор найденных ошибок: что и почему ломалось}

Раздел полезен тем, что показывает \emph{типичные} ошибки в подобном коде и
метод их поиска. Все случаи реальные, из истории разработки проекта.

\subsection{Выход за границы вектора в трекере}

\paragraph{Симптом.} Тест трекера завершается с \texttt{Segmentation fault}.

\paragraph{Как искали.} Отладчика \texttt{gdb} в среде не было. Сборка с
\texttt{-fsanitize=address} дала точное место: обращение к
\texttt{std::vector<bool>} по индексу за границей.

\paragraph{Причина.} Флаги <<трек был сопоставлен>> создавались по размеру
списка треков \emph{до} добавления новых. Затем цикл старения проходил по
\emph{новому} размеру, включающему только что добавленные треки.

\begin{lstlisting}
std::vector<bool> track_used(tracks_.size(), false);   // старый размер
// ... добавление новых треков: tracks_.push_back(...)
for (std::size_t t = 0; t < tracks_.size(); ++t) {     // новый размер! UB
    if (!track_used[t]) { ++tracks_[t].missed; }
}
\end{lstlisting}

\paragraph{Исправление.} Граница цикла привязана к размеру \emph{флагов},
которые описывают именно <<треки, существовавшие до этого кадра>>.

\paragraph{Мораль.} Индексная переменная должна ограничиваться тем
контейнером, по которому идёт индексация. \texttt{operator[]} у
\texttt{std::vector} не проверяет границы (в отличие от \texttt{at()}).

\subsection{\texttt{terminate} при остановке конвейера}

\paragraph{Симптом.} При выходе: \texttt{terminate called without an active
exception}.

\paragraph{Причина.} \texttt{join()} вызывался до снятия флага остановки.
Рабочие циклы крутились по флагу, которого никто не снимал; кроме того, при
некотором раскладе рабочий цикл завершался сам, и объект
\texttt{std::thread} оставался \emph{joinable} --- это и вызывает
\texttt{std::terminate} при разрушении.

\paragraph{Исправление.} Порядок строго: снять флаг, затем \texttt{join}.

\subsection{Трекер <<залипал>>: пропущенное уравнение ковариации}

Самая содержательная ошибка --- она математическая, а не программная.
Подробный разбор --- в разделе о фильтре Калмана: вызов
\texttt{kalman.predict()} был заменён на умножение матрицы перехода, из-за
чего ковариация ошибки не росла, усиление было близко к нулю, и трек не
следовал за объектом.

\paragraph{Почему не замечалось сразу.} На медленном движении (несколько
пикселей за кадр) IoU оставался выше порога, и трек не терялся. Ошибка
проявилась бы на реальной камере, где объекты смещаются заметно.

\paragraph{Мораль.} Тест должен воспроизводить \emph{напряжённый} режим
(в данном случае --- быстрое движение), иначе он не отличает правильную
реализацию от почти правильной.

\subsection{Нулевая скорость в HUD}

\paragraph{Симптом.} Счётчики кадров в секунду всегда показывали $0{,}0$.

\paragraph{Причина.} Метрики пересчитывались только внутри ветки обработки
полученного результата, а точка отсчёта обновлялась там же. В итоге знаменатель
оказывался нулевым или период не накапливался.

\paragraph{Мораль.} Периодические измерения нужно обновлять независимо от
того, пришли ли новые данные; иначе метрика <<замерзает>> именно в тот
момент, когда она важнее всего (поток данных встал).

\subsection{Отрисовка масок: полный кадр вместо рамки}

\paragraph{Симптом.} На модели, выдающей много объектов (несколько сотен),
обработка падает со 30 кадров/с до 6--7, а счётчик вытеснений растёт.

\paragraph{Что измеряли.} По HUD было видно \texttt{objects 300}, то есть
дело не в декодере. Замеры по частям:

\begin{center}
\begin{tabular}{lr}
\toprule
Операция (1280$\times$720, 300 объектов) & Время на кадр \\
\midrule
трекер & 1{,}7 мс \\
отрисовка без масок & 10{,}8 мс \\
отрисовка 300 масок (\emph{до} исправления) & 266{,}5 мс \\
отрисовка 300 масок (\emph{после}) & 17{,}7 мс \\
\bottomrule
\end{tabular}
\end{center}

\paragraph{Причина.} Маска хранится в размере \emph{всего} кадра, хотя
ненулевые пиксели занимают лишь рамку объекта. Код же делал
\texttt{setTo(color, mask)} и \texttt{findContours(mask)} по полному кадру.
Стоимость: $300 \times 1280 \times 720 \approx 2{,}8\cdot 10^{8}$ операций на
кадр --- при том что полезная площадь на два порядка меньше.

\paragraph{Исправление.} Обе операции выполняются над ROI --- прямоугольником
рамки. Контуры, найденные в ROI, затем сдвигаются на смещение рамки.

\begin{lstlisting}
const cv::Rect roi = instance.box & cv::Rect(0, 0, canvas.cols, canvas.rows);
cv::Mat layer_roi = layer(roi);
layer_roi.setTo(class_color(instance.class_id), instance.mask(roi));
\end{lstlisting}

\begin{note}
\texttt{cv::Mat::operator()} не копирует пиксели: ROI --- это новый заголовок
матрицы, ссылающийся на ту же память, но со своим шагом строки. Поэтому запись
в \texttt{layer\_roi} меняет исходный \texttt{layer}. Это идиома <<вид на
область>>, и её надо знать: изменение вида изменяет оригинал.
\end{note}

\paragraph{Мораль.} Сложность алгоритма нужно считать по \emph{реально
обрабатываемой} площади, а не по размеру контейнера. И измерять по частям:
без замеров легко <<оптимизировать>> не то место --- трекер, который был
подозреваемым номером один, занимал менее 2\% времени.

\subsection{Молчаливая потеря объектов при неожиданном выходе модели}

\paragraph{Симптом.} При нестандартной форме второго выхода модели на экране
нет ни объектов, ни сообщений об ошибке.

\paragraph{Причина.} Результат заполнялся \emph{только} внутри ветки
<<построить маски>>. Если второй выход не был картой прототипов, функция
возвращала пустой список.

\paragraph{Исправление.} Явная проверка формы выхода и \emph{деградация}:
при отсутствии прототипов отдаются хотя бы рамки, а в лог пишется
предупреждение.

\paragraph{Мораль.} Любая <<деградация>> должна быть видимой. Пустой
результат без сообщения --- худший вид ошибки: он выглядит как корректная
работа на кадре без объектов.

\section{Как расширять проект}

Разделение на интерфейсы выбрано так, чтобы типичные изменения затрагивали
один файл.

\subsection{Добавить новый источник кадров}

Достаточно реализовать \texttt{Source} и выбрать реализацию в \texttt{main}:

\begin{lstlisting}
class FileSource final : public Source {
public:
    bool open(std::string& error) override;
    bool grab(FramePtr& out, std::string& error) override;
    int width() const override;
    int height() const override;
};
\end{lstlisting}

Конвейер, трекер и отрисовка при этом не меняются: они работают с
\texttt{Source\&}.

\subsection{Поддержать другой формат выходов модели}

Добавляется новая ветка декодирования и новое значение параметра
\texttt{out\_layout}. Важно, что формат задаётся \emph{явно}: пользователь
обязан указать, какой у него экспорт, а программа печатает фактические формы
тензоров при загрузке --- расхождение видно сразу.

\subsection{Улучшить сопоставление в трекере}

Жадный алгоритм локализован в одной функции. Замены по возрастанию сложности:

\begin{enumerate}
  \item венгерский алгоритм (оптимальное назначение по суммарному IoU);
  \item каскадное сопоставление: сначала по IoU, затем по расстоянию между
        центрами;
  \item учёт класса и цвета/гистограмм;
  \item DeepSORT с эмбеддингами внешнего вида.
\end{enumerate}

Все они вставляются на место шага сопоставления и не затрагивают очередь,
потоки или отрисовку.

\subsection{Что менять не следует}

\begin{itemize}
  \item \textbf{Политику переполнения очереди} --- вытеснение старого
        является смыслом схемы реального времени.
  \item \textbf{Неизменяемость кадра} на границе потоков: именно она
        обеспечивает потокобезопасность без блокировок.
  \item \textbf{Явность форматов}: любой автоподбор скрывает причину ошибок.
\end{itemize}

\section{Итог: принципы, которые переносятся на другие проекты}

\begin{enumerate}
  \item \textbf{Разделяй стадии по времени отклика.} Захват (жёсткий темп)
        и обработка (переменное время) --- разные потоки.
  \item \textbf{Ограничивай буферы и выбирай политику переполнения осознанно.}
        Для реального времени правильная политика --- вытеснение старого.
  \item \textbf{Делай данные неизменяемыми на границе потоков.}
        \texttt{shared\_ptr<const T>} убирает гонки без блокировок.
  \item \textbf{Никогда не блокируй производителя.} Иначе камера встанет.
  \item \textbf{Задавай форматы явно.} Автодетект скрывает ошибки.
  \item \textbf{Пиши тест на каждую найденную ошибку.} Тест на движение в
        трекере и тест баланса очереди появились именно так.
  \item \textbf{Переводи исключения в коды возврата на границе потоков.}
        Иначе \texttt{std::terminate}.
  \item \textbf{Измеряй задержку, а не только пропускную способность.}
        Система может быть быстрой в среднем и непригодной в реальном времени.
\end{enumerate}

\appendix

\section{Сводка формул}

\begin{align}
  \text{letterbox:}\quad & s = \min(S/W, S/H), \quad x = su + p_x \\
  \text{обратно:}\quad & u = (x - p_x)/s \\
  \text{IoU:}\quad & \frac{|A \cap B|}{|A| + |B| - |A \cap B|} \\
  \text{сигмоида:}\quad & \sigma(z) = \frac{1}{1+e^{-z}} \\
  \text{порог маски:}\quad & \sigma(z) > t \iff z > \ln\frac{t}{1-t} \\
  \text{Калман, предсказание:}\quad & \hat{x}^- = F\hat{x}, \quad P^- = FPF^\top + Q \\
  \text{Калман, коррекция:}\quad & K = P^-H^\top(HP^-H^\top + R)^{-1} \\
  & \hat{x} = \hat{x}^- + K(z - H\hat{x}^-), \quad P = (I-KH)P^- \\
  \text{Precision/Recall:}\quad & \frac{TP}{TP+FP}, \quad \frac{TP}{TP+FN} \\
  \text{mAP:}\quad & \frac{1}{n_c}\sum_c \mathrm{AP}_c
\end{align}

\section{Словарь терминов}

\begin{description}[leftmargin=3.2cm,style=nextline]
  \item[blob] непрерывный массив чисел, здесь --- вход сети в формате NCHW.
  \item[letterbox] вписывание изображения в квадрат с сохранением пропорций и
        полями.
  \item[NMS] non-maximum suppression, подавление перекрывающихся детекций.
  \item[IoU] intersection over union, мера перекрытия двух областей.
  \item[anchor] позиция на карте признаков, для которой сеть предсказывает
        объект.
  \item[прототип маски] одна из 32 базовых масок, линейная комбинация которых
        даёт маску объекта.
  \item[RAII] привязка владения ресурсом к времени жизни объекта.
  \item[pimpl] идиома сокрытия деталей реализации за неполным типом.
  \item[вытеснение] удаление самого старого элемента при переполнении очереди.
  \item[rpath] путь поиска динамических библиотек, вшитый в исполняемый файл.
\end{description}

\end{document}
